\chapter{Action receipts and exact closure}\label{ch:receipts}

Successive receipts form a physical history because each event starts where the preceding one ends. Their values therefore telescope. This exact closure is the foundation for reconciling local EBU records, detecting duplicate claims and preserving the contribution of necessary actions with either sign.


\section{What a physical receipt records}
A shop receipt records a commercial exchange. An EBU action receipt names a
physical predecessor, an accepted transformation, its successor and the
version of the potential used to compare them. The two kinds of receipt may
coexist, but one does not validate the other. A payment cannot prove that a
reservoir recovered, and a positive EBU does not establish who was paid.

The preceding factor chapters showed why a complete change can often be evaluated
from affected factors alone. We now connect several correctly evaluated
changes. The central fact is simple: the state after one action is the state
before the next. Its potential appears once with each sign, and cancels.

\section{Follow the actual chain}
Consider a sequence of authoritative states under one fixed potential,
\begin{equation}
 x_0\longrightarrow x_1\longrightarrow\cdots\longrightarrow x_m.
\end{equation}
For each complete transition define
\begin{equation}
 E_r=V(x_{r-1})-V(x_r).
\end{equation}
The arrows are physical or modeled events, not an arbitrary sort order for
files. If several actions occur together, one arrow can be their joint
transition, retaining the identities of its children. A later allocation of
that group's total is another question.

Summing the consecutive differences gives
\begin{align}
 \sum_{r=1}^m E_r
 &=V(x_0)-V(x_1)+V(x_1)-V(x_2)+\cdots
   +V(x_{m-1})-V(x_m)\\
 &=V(x_0)-V(x_m).
\end{align}
Equivalently,
\begin{equation}\label{eq:receipt-closure}
 \sum_r E_r+V(x_m)-V(x_0)=0.
\end{equation}
This is exact telescoping for any state function $V$. Gaussian geometry is
not needed. Neither is an actor willing to improve the state. The
intermediate states still matter: without them we cannot test whether each
action was physical, whether its recorded value was correct or whether an
event is missing.

\begin{lawbox}{What closure proves}
Correctly linked, complete transitions under one fixed potential have a
total EBU equal to their endpoint potential reduction. The identity gives the account an exact reconciliation rule. Ownership, service and recovery retain their separately declared meanings.
\end{lawbox}

\section{One sequence, not three fresh beginnings}
Use the two-cell potential with reference $(10,10)$ and scales $(2,2)$.
Three successive lossless transfers carry the state through
\[
 (18,2)\to(16,4)\to(13,7)\to(12,8),
 \qquad V:16\to9\to\frac94\to1.
\]
The action values are $7$, $27/4$ and $5/4$, which add to $15=16-1$.
The middle transfer starts from $(16,4)$, not from $(18,2)$. Reusing the
initial slope for it would quote a different imagined action, not the
second event in this sequence.

\diagram{sequence}{Three consecutive finite differences telescope to fifteen. Every action uses the preceding endpoint as its baseline.}{fig:sequence}

Suppose a record assigned the first and third lines to one actor and the
second to another. Their historical contributions would be $33/4$ and
$27/4$, adding to the same $15$. This partition uses an explicit owner rule;
the potential did not choose the owners. Individual accounts can remain
signed. There is no physical theorem requiring a negative contribution to
be prohibited or prefunded.

\section{Why the negative line must remain}
Let one action take potential from $1$ to $4$ and another take it from $4$ to $0$. Their EBU values are $-3$ and $+4$. The total is $+1$, exactly the endpoint reduction. Keeping only the positive line would report four units of improvement where the complete history contains one.

An institution can support the first action when it supplies necessary service. Support belongs in its allocation rules; the sign belongs in the physical history. Maintaining both gives the institution a more useful record: it can see what restoration achieved and what service required.

The same principle covers autonomous feedback. If a registered action starts a response that later overshoots and settles, successive observations can explain the path, but they are not automatically new actor events. A no-action background can have nonzero physical change. Its line must remain distinguishable from the incremental action comparison.

Exact closure is therefore socially useful as well as algebraically neat. It makes selective reporting harder, allows small local records to be reconciled into a common total, and keeps necessary burdens visible so they can be reduced through better design rather than hidden through accounting.

\section{External events need their own line}
Now let an external event take $x_t$ to $x^{\mathrm{b}}_t$, followed by an actor action
from $x^{\mathrm{b}}_t$ to $x_{t+1}$. Define the external potential change by
\[
 D_{{\mathrm{ext}},t}=V(x^{\mathrm{b}}_t)-V(x_t).
\]
The actor's EBU is $E_t=V(x^{\mathrm{b}}_t)-V(x_{t+1})$, so
\begin{equation}\label{eq:driven-closure}
 \sum_t E_t=V(x_0)-V(x_N)+\sum_tD_{{\mathrm{ext}},t}.
\end{equation}
An external worsening can create deviation that a later actor reduces.
An external improvement can remove deviation without actor contribution.
Both signs belong in the record. For example, if an external redistribution
moves potential from $0$ to $4$ and an actor then reduces it to $1$, the
actor records $3$, not the $-1$ obtained by comparing the wrong endpoints.

If the potential itself changes between events, the calculation needs an
explicit field-change line too. Comparing values produced by two different
rulers as though they were one fixed $V$ breaks Equation~\ref{eq:receipt-closure}.
We return to that harder question after the fixed-field theory is complete.

\section{Balance is not carrier conservation}
In a lossless transfer, physical stock is conserved while $V$ changes.
States $(18,2)$ and $(14,6)$ each contain twenty stock units; their
potentials are sixteen and four. A twelve-EBU reduction did not create or
destroy twelve litres. Equation~\ref{eq:receipt-closure} is a reconciliation
of values, not a conservation law for matter or energy.

An auditor therefore checks both ledgers. First follow the physical states,
routes, carriers and boundaries. Then verify potential versions, event
identities and signed EBU lines. A residual can expose an omitted event,
duplicate receipt, stale baseline or arithmetic error; it cannot by itself
reconstruct the missing physical cause. The next chapter uses this discipline
to examine closed loops and action splitting.
